\documentclass[11pt, a4paper]{article}
\usepackage[utf8]{inputenc}
\usepackage{amsmath, amssymb}
\usepackage{graphicx}
\usepackage{hyperref}
\usepackage{geometry}
\geometry{margin=1in}
\usepackage{microtype}

\title{EE 325 Digital Signal Processing \\ Assignments Report}
\author{PERERA J.D.T. \\ E/21/291}
\date{\today}

\begin{document}
\maketitle
\tableofcontents
\newpage

% Assignment 1
\section{Assignment 1: Sampling and Aliasing}
\subsection{Objective}
Demonstrate the effects of sampling a continuous-time signal at frequencies below the Nyquist rate (Aliasing).

\subsection{Problem Description}
\textbf{Continuous Signal:}
A sine wave of frequency $f = 1000$ Hz is used:
$$x_c(t) = \sin(2\pi \cdot 1000 t)$$
\textbf{Nyquist Rate:}
The minimum sampling rate required to avoid aliasing is $2f_{max} = 2 \times 1000 = 2000$ Hz.

\subsection{Simulation Analysis}
The Python script sampled this signal at two different frequencies:

\subsubsection{Case 1: Sampling at $f_s = 600$ Hz}
\begin{itemize}
    \item \textbf{Observation:} Since $f_s = 600$ Hz $< 2000$ Hz, the system is \textbf{under-sampled}.
    \item \textbf{Result:} Aliasing occurs. The high-frequency original signal ($1000$ Hz) will appear as a lower frequency alias in the discrete domain.
    \item \textbf{Alias Frequency:} $|1000 - k \cdot 600|$. For $k=2$, $|1000 - 1200| = 200$ Hz (folded).
\end{itemize}

\subsubsection{Case 2: Sampling at $f_s = 1500$ Hz}
\begin{itemize}
    \item \textbf{Observation:} Since $f_s = 1500$ Hz $< 2000$ Hz, the system is also \textbf{under-sampled}.
    \item \textbf{Result:} Aliasing occurs.
    \item \textbf{Alias Frequency:} $|1000 - 1500| = 500$ Hz.
\end{itemize}

\subsection{Conclusion}
The simulation visually demonstrates that sampling at a rate lower than the Nyquist rate results in a distorted or ``aliased'' signal that does not faithfully represent the original frequency. To correctly reconstruct the 1000 Hz sine wave, a sampling rate of at least 2000 Hz is required.

\newpage
% Assignment 2
\section{Assignment 2: Time Domain Analysis of LTI Systems}
\subsection{Objective}
Analyze the time-domain response of a Linear Time-Invariant (LTI) system using difference equations in Python.

\subsection{System Description}
The system is defined by the first-order difference equation:
$$y[n] = 0.5 y[n-1] + x[n]$$
Rearranging terms, this corresponds to the system:
$$y[n] - 0.5 y[n-1] = x[n]$$
Impulse Response expectation: $h[n] = (0.5)^n u[n]$.

\subsection{Simulation Results}
\subsubsection{1. Impulse Response ($h[n]$)}
\begin{itemize}
    \item \textbf{Input:} $x[n] = \delta[n]$ (Unit Impulse).
    \item \textbf{Method:} Iteratively inferred $y[n]$ with $y[-1]=0$.
    \item \textbf{Result:}
    \begin{itemize}
        \item $n=0: y[0] = 0.5(0) + 1 = 1$
        \item $n=1: y[1] = 0.5(1) + 0 = 0.5$
        \item $n=2: y[2] = 0.5(0.5) + 0 = 0.25$
    \end{itemize}
    \item \textbf{Verification:} Follows the theoretical decay $h[n] = (0.5)^n$.
\end{itemize}

\subsubsection{2. Step Response ($s[n]$)}
\begin{itemize}
    \item \textbf{Input:} $x[n] = u[n]$ (Unit Step).
    \item \textbf{Result:} Accumulates the impulse response.
    \begin{itemize}
        \item $n=0: 1$
        \item $n=1: 1 + 0.5 = 1.5$
        \item $n=2: 1.5 + 0.25 = 1.75$
    \end{itemize}
    \item \textbf{Steady State:} As $n \to \infty$, $s[n] \to \frac{1}{1 - 0.5} = 2$.
\end{itemize}

\subsection{Conclusion}
The simulation successfully implements the recursive difference equation and visualizes the system's characteristic exponential decay (Stable IIR system).

\newpage
% Assignment 3
\section{Assignment 3: Discrete Time Convolution}
\subsection{Objective}
Verify the calculation of linear convolution for discrete-time sequences using numerical tools.

\subsection{Q1: Linear Convolution}
\textbf{Sequences:}
\begin{itemize}
    \item $x[n] = [1, 2, 1]$
    \item $h[n] = [1, -1, 2]$
\end{itemize}

\textbf{Convolution $y[n] = x[n] * h[n]$:}
\begin{itemize}
    \item Length of $y$: $L_x + L_h - 1 = 3 + 3 - 1 = 5$.
    \item Result:
    \begin{itemize}
        \item $y[0] = 1$
        \item $y[1] = 1$
        \item $y[2] = 1$
        \item $y[3] = 3$
        \item $y[4] = 2$
    \end{itemize}
    \item \textbf{Verified Output:} $[1, 1, 1, 3, 2]$
\end{itemize}

\subsection{Q2: Cascade Connection of Systems}
\textbf{Systems:}
\begin{itemize}
    \item $h_1[n] = [1, 1]$
    \item $h_2[n] = [1, -1, 1]$
\end{itemize}

\textbf{Equivalent Impulse Response $h_{eq}[n]$:}
For cascade systems, $h_{eq}[n] = h_1[n] * h_2[n]$.
\begin{itemize}
    \item Length: $2 + 3 - 1 = 4$.
    \item Calculation:
    \begin{itemize}
        \item $h_{eq}[0] = 1$
        \item $h_{eq}[1] = 0$
        \item $h_{eq}[2] = 0$
        \item $h_{eq}[3] = 1$
    \end{itemize}
    \item \textbf{Verified Output:} $[1, 0, 0, 1]$
\end{itemize}

\subsection{Conclusion}
The computational results confirm that numerical convolution matches the analytical expectations for both single-stage and cascade LTI systems.

\newpage
% Assignment 4
\section{Assignment 4: z-Transform and Inverse z-Transform}

\subsection{1. Forward z-Transform}
\textbf{Question:} Find the z-transform $X(z)$ and region of convergence (ROC) of $x[n] = 0.5^n u[n]$.

\textbf{Solution:}
The definition of the z-transform is:
$$X(z) = \sum_{n=-\infty}^{\infty} x[n] z^{-n}$$

Substituting $x[n]$:
$$X(z) = \sum_{n=0}^{\infty} 0.5^n z^{-n} = \sum_{n=0}^{\infty} (0.5 z^{-1})^n$$

This is a geometric series of the form $\sum_{n=0}^{\infty} r^n = \frac{1}{1-r}$, which converges if $|r| < 1$.
Here, $r = 0.5z^{-1}$.

Convergence condition:
$$|0.5 z^{-1}| < 1 \implies \frac{0.5}{|z|} < 1 \implies |z| > 0.5$$

Result:
$$X(z) = \frac{1}{1 - 0.5z^{-1}} = \frac{z}{z - 0.5}$$

\textbf{ROC:} $|z| > 0.5$

\subsection{2. Inverse z-Transform by Partial Fractions}
\textbf{Question:} Given $X(z) = \frac{2z}{(z-1)(z+0.5)}$ and assuming valid causal $x[n]$, find $x[n]$.

\textbf{(a) Partial Fraction Expansion:}
We expand $\frac{X(z)}{z}$ to easily handle the roots:
$$\frac{X(z)}{z} = \frac{2}{(z-1)(z+0.5)} = \frac{A}{z-1} + \frac{B}{z+0.5}$$

finding A:
$$A = \left. \frac{2}{z+0.5} \right|_{z=1} = \frac{2}{1.5} = \frac{4}{3}$$

finding B:
$$B = \left. \frac{2}{z-1} \right|_{z=-0.5} = \frac{2}{-1.5} = -\frac{4}{3}$$

So,
$$\frac{X(z)}{z} = \frac{4/3}{z-1} - \frac{4/3}{z+0.5}$$
$$X(z) = \frac{4}{3} \frac{z}{z-1} - \frac{4}{3} \frac{z}{z+0.5}$$

\textbf{(b) Time-domain sequence x[n]:}
Using standard transform pairs for causal sequences ($|z| > 1$):
$$\frac{z}{z-a} \longleftrightarrow a^n u[n]$$

Term 1: $\frac{4}{3} (1)^n u[n]$ \\
Term 2: $-\frac{4}{3} (-0.5)^n u[n]$

$$x[n] = \frac{4}{3} u[n] - \frac{4}{3} (-0.5)^n u[n]$$
$$x[n] = \frac{4}{3} [1 - (-0.5)^n] u[n]$$

\textbf{(c) Sketch and Decay:}
\begin{itemize}
    \item For $n=0$: $x[0] = 4/3(1 - 1) = 0$
    \item For $n=1$: $x[1] = 4/3(1 - (-0.5)) = 4/3(1.5) = 2$
    \item For $n \to \infty$: $(-0.5)^n \to 0$, so $x[n] \to 4/3 \approx 1.33$
\end{itemize}

The system has a transient response due to the pole at $-0.5$ (alternating decay) and a steady-state response due to the pole at $1$ (step).

\subsection{3. Python Verification}
The symbolic result matches the analytical result perfectly. Programmatically derived:
$$x[n] = 1.333 \cdot 1^n - 1.333 \cdot (-0.5)^n$$

\newpage
% Assignment 5
\section{Assignment 5: System Characterization via z-Transform}

\subsection{1. Transfer Function from Difference Equation}
\textbf{System:}
$$y[n] - 1.1y[n-1] + 0.3y[n-2] = x[n] + 0.5x[n-1]$$

\textbf{(a) Derive H(z):}
Taking the z-transform of both sides:
$$Y(z) - 1.1z^{-1}Y(z) + 0.3z^{-2}Y(z) = X(z) + 0.5z^{-1}X(z)$$
$$Y(z) [1 - 1.1z^{-1} + 0.3z^{-2}] = X(z) [1 + 0.5z^{-1}]$$

The transfer function $H(z) = \frac{Y(z)}{X(z)}$ is:
$$H(z) = \frac{1 + 0.5z^{-1}}{1 - 1.1z^{-1} + 0.3z^{-2}}$$

\textbf{(b) Coefficients:}
Expressing in the form $H(z) = \frac{b_0 + b_1 z^{-1} + \dots}{1 + a_1 z^{-1} + a_2 z^{-2}}$:
$$H(z) = \frac{1 + 0.5z^{-1}}{1 - 1.1z^{-1} + 0.3z^{-2}}$$

Coefficients:
\begin{itemize}
    \item Numerator $\{b_k\}$: $b_0 = 1, b_1 = 0.5$
    \item Denominator $\{a_k\}$: $a_0=1, a_1 = -1.1, a_2 = 0.3$
\end{itemize}

\subsection{2. Pole-Zero Analysis and Stability}

\textbf{(a) Poles and Zeros:}
Zeros (roots of numerator $1 + 0.5z^{-1}$):
$$1 + 0.5z^{-1} = 0 \implies z = -0.5$$
Zero at $z = -0.5$.

Poles (roots of denominator $1 - 1.1z^{-1} + 0.3z^{-2}$):
Multiply by $z^2$: $z^2 - 1.1z + 0.3 = 0$.
Roots using quadratic formula:
$$z = \frac{1.1 \pm \sqrt{1.1^2 - 4(1)(0.3)}}{2} = \frac{1.1 \pm \sqrt{1.21 - 1.2}}{2} = \frac{1.1 \pm \sqrt{0.01}}{2}$$
$$z = \frac{1.1 \pm 0.1}{2}$$
$p_1 = \frac{1.2}{2} = 0.6$\\
$p_2 = \frac{1.0}{2} = 0.5$

Poles at $z = 0.6$ and $z = 0.5$.

\textbf{(b) Stability:}
A causal LTI system is BIBO stable if all poles lie \textbf{strictly inside the unit circle} ($|p| < 1$).
Here, poles are $0.5$ and $0.6$.
Since $|0.5| < 1$ and $|0.6| < 1$, the system is \textbf{Stable}.

\textbf{(c) System Type:}
The system has poles (denominator order $> 0$), meaning $y[n]$ depends on past outputs. Thus, it has an impulse response of infinite duration. The system is \textbf{IIR (Infinite Impulse Response)}.

\subsection{3. Frequency Response}
\textbf{Gain Analysis:}
\begin{itemize}
    \item \textbf{DC Gain ($\omega = 0, z = 1$):} \\
    $$H(e^{j0}) = H(1) = \frac{1 + 0.5}{1 - 1.1 + 0.3} = \frac{1.5}{0.2} = 7.5$$ \\
    Magnitude: $20 \log_{10}(7.5) \approx 17.5$ dB.
    
    \item \textbf{Nyquist Gain ($\omega = \pi, z = -1$):} \\
    $$H(e^{j\pi}) = H(-1) = \frac{1 - 0.5}{1 + 1.1 + 0.3} = \frac{0.5}{2.4} \approx 0.208$$ \\
    Magnitude: $20 \log_{10}(0.208) \approx -13.6$ dB.
\end{itemize}

\textbf{Pole-Zero Influence:}
\begin{itemize}
    \item The poles are at $z=0.5, 0.6$. These are positive real poles, closest to the unit circle at $\omega=0$ ($z=1$).
    \item This proximity to $\omega=0$ causes a \textbf{rise in magnitude at low frequencies} (Low Pass behavior), which explains the high DC gain.
    \item The zero at $z=-0.5$ is on the negative real axis ($\omega=\pi$). This suppresses high frequencies, contributing to the low magnitude at $\pi$.
\end{itemize}

\newpage
% Assignment 6
\section{Assignment 6: Realization of Discrete-Time Systems}

\subsection{1. Direct Form I Realization}
\textbf{Difference Equation:}
$$y[n] - 1.1y[n-1] + 0.3y[n-2] = x[n] + 0.5x[n-1]$$
$$y[n] = 1.1y[n-1] - 0.3y[n-2] + x[n] + 0.5x[n-1]$$

\textbf{Block Diagram (Direct Form I):}
Direct Form I implements the feedforward part ($x[n]$ terms) and feedback part ($y[n]$ terms) separately, connected in cascade.

Specifically for $y[n] = 1.1y[n-1] - 0.3y[n-2] + ...$:
\begin{itemize}
    \item Tap from $y[n-1]$ is multiplied by $1.1$.
    \item Tap from $y[n-2]$ is multiplied by $-0.3$.
\end{itemize}

\textbf{Delay Elements:} 3 total ($1$ for $x$, $2$ for $y$). Note: This is non-canonic.

\subsection{2. Direct Form II Realization}
\textbf{(a) Direct Form II - Canonic:}
Direct Form II shares the delay elements between the feedforward and feedback paths. We introduce an intermediate variable $w[n]$.
$$w[n] = x[n] + 1.1w[n-1] - 0.3w[n-2]$$
$$y[n] = w[n] + 0.5w[n-1]$$

\textbf{(b) Memory Comparison:}
\begin{itemize}
    \item \textbf{Direct Form I:} Requires $M + N = 1 + 2 = 3$ delay elements ($x[n-1]$, $y[n-1]$, $y[n-2]$).
    \item \textbf{Direct Form II:} Requires $\max(M, N) = \max(1, 2) = 2$ delay elements ($w[n-1]$, $w[n-2]$).
    \item \textbf{Conclusion:} Direct Form II is more memory efficient (Canonical).
\end{itemize}

\textbf{(c) Numerical Differences:}
Direct Form II is more sensitive to quantization effects (overflow/round-off) at the internal nodes ($w[n]$) compared to Direct Form I, especially for high-Q poles, because the poles solely determine the magnitude of the intermediate variable $w[n]$ before zeros can attenuate it.

\newpage
% Assignment 7
\section{Assignment 7: Stability of Discrete-Time Systems}

\subsection{1. BIBO Stability from Impulse Response}
A linear time-invariant (LTI) system is BIBO stable if and only if the impulse response is absolutely summable:
$$ \sum_{n=-\infty}^{\infty} |h[n]| < \infty $$

\subsubsection{(a) $h_1[n] = 2^n u[n]$}
$$ \sum_{n=-\infty}^{\infty} |2^n u[n]| = \sum_{n=0}^{\infty} 2^n = 1 + 2 + 4 + \dots = \infty $$
\textbf{Conclusion:} Unstable. The sum diverges.

\subsubsection{(b) $h_2[n] = (-0.5)^n u[n]$}
$$ \sum_{n=-\infty}^{\infty} |(-0.5)^n u[n]| = \sum_{n=0}^{\infty} |(-0.5)|^n = \sum_{n=0}^{\infty} (0.5)^n $$
Using geometric series sum $\sum_{n=0}^{\infty} r^n = \frac{1}{1-r}$ for $|r| < 1$:
$$ \sum_{n=0}^{\infty} (0.5)^n = \frac{1}{1-0.5} = \frac{1}{0.5} = 2 < \infty $$
\textbf{Conclusion:} Stable. The sum is finite.

\subsubsection{(c) $h_3[n] = u[n]$}
$$ \sum_{n=-\infty}^{\infty} |u[n]| = \sum_{n=0}^{\infty} 1 = 1 + 1 + 1 + \dots = \infty $$
\textbf{Conclusion:} Unstable. The sum diverges.

\subsection{2. Jury Stability Test}

Characteristic polynomial: $p(z) = z^2 - 1.7z + 0.6$. \\
Coefficients: $a_2 = 1, a_1 = -1.7, a_0 = 0.6$.

\subsubsection{(a) Jury Table and Conditions}
For a 2nd order system ($N=2$), the necessary and sufficient conditions for stability are:
\begin{enumerate}
    \item $P(1) > 0$
    \item $(-1)^2 P(-1) > 0 \Rightarrow P(-1) > 0$
    \item $|a_0| < a_2$ (assuming $a_2 > 0$)
\end{enumerate}

\textbf{Checking Conditions:}
\begin{enumerate}
    \item $P(1) = (1)^2 - 1.7(1) + 0.6 = 1 - 1.7 + 0.6 = -0.1$ \\
    Condition $P(1) > 0$ is \textbf{FALSE}.
    \item $P(-1) = (-1)^2 - 1.7(-1) + 0.6 = 1 + 1.7 + 0.6 = 3.3$ \\
    Condition $P(-1) > 0$ is \textbf{TRUE}.
    \item $|a_0| < a_2 \Rightarrow |0.6| < 1$ \\
    Condition is \textbf{TRUE}.
\end{enumerate}

\textbf{Conclusion:} Since the first condition fails ($P(1) = -0.1 \ngtr 0$), the system is \textbf{UNSTABLE}.

\subsubsection{(b) Verification with Roots}
Roots of $z^2 - 1.7z + 0.6 = 0$:
$$ z = \frac{1.7 \pm \sqrt{(-1.7)^2 - 4(1)(0.6)}}{2} = \frac{1.7 \pm \sqrt{2.89 - 2.4}}{2} $$
$$ z = \frac{1.7 \pm \sqrt{0.49}}{2} = \frac{1.7 \pm 0.7}{2} $$
$$ z_1 = \frac{2.4}{2} = 1.2, \quad z_2 = \frac{1.0}{2} = 0.5 $$
Only stable if all poles satisfy $|z| < 1$. \\
Here, $|z_1| = 1.2 > 1$. \\
\textbf{Conclusion:} The system is \textbf{UNSTABLE}, consistent with the Jury test.

\subsection{3. System Behavior Analysis}
\begin{itemize}
    \item \textbf{System A:} $y[n] - 0.5y[n-1] = x[n]$
    \begin{itemize}
        \item Characteristic equation: $z - 0.5 = 0 \Rightarrow z = 0.5$.
        \item Since $|z| = 0.5 < 1$, the pole is inside the unit circle.
        \item \textbf{Expectation:} Stable system, bounded output for bounded input.
    \end{itemize}
    \item \textbf{System B:} $y[n] - 1.2y[n-1] = x[n]$
    \begin{itemize}
        \item Characteristic equation: $z - 1.2 = 0 \Rightarrow z = 1.2$.
        \item Since $|z| = 1.2 > 1$, the pole is outside the unit circle.
        \item \textbf{Expectation:} Unstable system, unbounded output (exponential growth) for bounded input.
    \end{itemize}
\end{itemize}

\newpage
% Assignment 8
\section{Assignment 8: Frequency-Domain Analysis – DTFT and DFS}

\subsection{1. DTFT of a Finite-Duration Signal}
The signal is a rectangular pulse of length $N$:
$$ x[n] = \begin{cases} 1 & 0 \le n \le N-1 \\ 0 & \text{otherwise} \end{cases} $$

\subsubsection{(a) Closed Form Derivation}
The DTFT is defined as:
$$ X(e^{j\omega}) = \sum_{n=-\infty}^{\infty} x[n]e^{-j\omega n} = \sum_{n=0}^{N-1} e^{-j\omega n} $$
This is a geometric series with ratio $r = e^{-j\omega}$ and first term $a=1$:
$$ X(e^{j\omega}) = \frac{1 - (e^{-j\omega})^N}{1 - e^{-j\omega}} = \frac{1 - e^{-j\omega N}}{1 - e^{-j\omega}} $$

\subsubsection{(b) Form involving sine functions and zero locations}
To express in terms of sine, we pull out half-angle exponential factors:
$$ X(e^{j\omega}) = \frac{e^{-j\omega N/2} (e^{j\omega N/2} - e^{-j\omega N/2})}{e^{-j\omega/2} (e^{j\omega/2} - e^{-j\omega/2})} $$
Using the Euler identity $\sin(\theta) = \frac{e^{j\theta} - e^{-j\theta}}{2j}$:
$$ X(e^{j\omega}) = \frac{e^{-j\omega N/2} \cdot 2j \sin(\omega N/2)}{e^{-j\omega/2} \cdot 2j \sin(\omega/2)} = e^{-j\omega(N-1)/2} \frac{\sin(\omega N/2)}{\sin(\omega/2)} $$

\textbf{Zero Locations:}
Zeros occur when the numerator is zero but the denominator is not.
$$ \sin(\omega N/2) = 0 \Rightarrow \frac{\omega N}{2} = k\pi \Rightarrow \omega = \frac{2\pi k}{N}, \quad k \in \mathbb{Z} $$
Exceptions: When $\sin(\omega/2) = 0$ (i.e., $\omega = 2\pi m$), the value is $N$ (using L'Hopital's rule), so these are not zeros.
Thus, zeros are at $\omega = \frac{2\pi k}{N}$ for $k \ne mN$.

For $N=8$, zeros in $[-\pi, \pi]$ are at:
$$ \omega = \pm \frac{2\pi}{8}, \pm \frac{4\pi}{8}, \pm \frac{6\pi}{8} \Rightarrow \pm \frac{\pi}{4}, \pm \frac{\pi}{2}, \pm \frac{3\pi}{4} $$

\subsubsection{(c) Main-lobe and Side-lobe Structure}
\begin{itemize}
    \item \textbf{Main Lobe:} Centered at $\omega = 0$. Extends from the first negative zero ($-\frac{2\pi}{N}$) to the first positive zero ($+\frac{2\pi}{N}$).
    \begin{itemize}
        \item Width = $\frac{4\pi}{N}$.
        \item Peak amplitude = $N$ at $\omega = 0$.
    \end{itemize}
    \item \textbf{Side Lobes:} Located between subsequent zeros. The amplitude of side lobes decreases as $\omega$ moves away from 0.
    \begin{itemize}
        \item The width of each side lobe is $\frac{2\pi}{N}$.
    \end{itemize}
\end{itemize}

\subsection{2. Discrete Fourier Series (DFS)}
Periodic signal with $N=4$:
$$ x[n] = \{1, 2, 3, 4\} \quad \text{for } n=0,1,2,3 $$

\subsubsection{(a) Compute DFS Coefficients}
The DFS definition is:
$$ \tilde{X}[k] = \sum_{n=0}^{N-1} \tilde{x}[n] e^{-j\frac{2\pi}{N}kn} = \sum_{n=0}^{3} x[n] e^{-j\frac{\pi}{2}kn} $$

\textbf{For k=0:}
$$ X[0] = \sum_{n=0}^{3} x[n] = 1 + 2 + 3 + 4 = 10 $$

\textbf{For k=1:}
$$ X[1] = 1(e^0) + 2(e^{-j\pi/2}) + 3(e^{-j\pi}) + 4(e^{-j3\pi/2}) $$
$$ X[1] = 1 + 2(-j) + 3(-1) + 4(j) = 1 - 3 + j(4-2) = -2 + 2j $$

\textbf{For k=2:}
$$ X[2] = 1(e^0) + 2(e^{-j\pi}) + 3(e^{-j2\pi}) + 4(e^{-j3\pi}) $$
$$ X[2] = 1 + 2(-1) + 3(1) + 4(-1) = 1 - 2 + 3 - 4 = -2 $$

\textbf{For k=3:}
$$ X[3] = 1(e^0) + 2(e^{-j3\pi/2}) + 3(e^{-j3\pi}) + 4(e^{-j9\pi/2}) $$
Note $e^{-j3\pi/2} = j$, $e^{-j3\pi} = -1$, $e^{-j9\pi/2} = e^{-j\pi/2} = -j$.
$$ X[3] = 1(1) + 2(j) + 3(-1) + 4(-j) = 1 - 3 + j(2-4) = -2 - 2j $$

\textbf{Result:}
$$ X[k] = \{10, -2+2j, -2, -2-2j\} $$

\subsubsection{(b) Real vs Complex and Symmetry}
\begin{itemize}
    \item $X[0]$ and $X[2]$ are real.
    \item $X[1]$ and $X[3]$ are complex.
    \item \textbf{Reason:} For a real-valued sequence $x[n]$, the DFT/DFS is conjugate symmetric:
    $$ X[k] = X^*[N-k] $$
    \begin{itemize}
        \item $X[1] = -2+2j$
        \item $X[4-1] = X[3] = -2-2j$
        \item indeed, $(-2+2j)^* = -2-2j$.
        \item $X[0]$ and $X[N/2]=X[2]$ are their own conjugates, so they must be real.
    \end{itemize}
\end{itemize}

\end{document}
